\section{Onde aprofundar}

\subsection*{Livros e teses}
\small
\begin{longtable}{@{}p{5.2cm} p{7cm}@{}}
\toprule
Fonte & Para quê \\
\midrule \endhead
\citetitle{livro:stenman} (Stenman, aberto, CC BY 4.0) & A BEAM por
  dentro: processos, escalonamento por reductions, GC, carga de código \\
\citetitle{paper:armstrong-thesis-reliable-systems} (Armstrong, KTH,
  2003) & A filosofia do Erlang/OTP; leia os capítulos~2 e~4 \\
\citetitle{livro:cesarini2016} (Cesarini e Vinoski, 2016) & Supervisão,
  \emph{behaviours} e confiabilidade em sistemas OTP reais \\
\citetitle{paper:history-of-erlang} (Armstrong, HOPL~III, 2007) & A
  história do Erlang contada pelo criador \\
\bottomrule
\end{longtable}
\normalsize

\subsection*{Documentação}
\small
\begin{longtable}{@{}p{5.2cm} p{7cm}@{}}
\toprule
Fonte & Para quê \\
\midrule \endhead
\citetitle{doc:erlang-processes}, \citetitle{doc:erlang-eff-processes}
  & Processos, sinais, links e custo de memória \\
\citetitle{doc:erlang-gc} & O coletor por processo e os binários
  grandes \\
\citetitle{doc:erlang-supervisor} & Estratégias de reinício e limite de
  intensidade \\
\citetitle{doc:erlang-code-loading} & As duas versões de um módulo \\
\citetitle{doc:erlang-beamasm} & Como funciona o JIT da BEAM \\
\citetitle{doc:graalvm-compiler}, \citetitle{doc:truffle} & O Graal como
  JIT e como base do Truffle \\
\citetitle{doc:jep505} & Concorrência estruturada no Java, inspirada em
  parte nos supervisores do Erlang \\
\bottomrule
\end{longtable}
\normalsize

\subsection*{Papers}
\textcite{paper:truffle-partial-evaluation} e
\textcite{paper:native-image-initialize-once} explicam, com dados, os
dois modos da GraalVM que este capítulo resumiu.
